% !TEX root = ../../vzor.tex
%% Ukázka: bloková schémata, UML diagram tříd, schéma databáze
%% Vkládá ji soubor vzor/kapitoly/prakticka.tex; ve vzor.pdf je pod nadpisem uvedeno „Zdroj ukázky".
%% Vyžaduje balíček tikz – v mprace.tex odkomentujte připravený řádek s \usepackage.

\section{Schémata programů}
\zdrojukazky{vzor/ukazky/schemata.tex}

Schémata jsou kreslená v \texttt{tikz} -- v souboru \texttt{mprace.tex}
odkomentujte řádek s \verb|\usepackage{tikz}|.

Dobrý popis složitějšího programu neobsahuje jen výpis kódu, ale také schéma struktury.
Schéma čtenáři rychle ukáže \emph{co} program dělá a \emph{jak} jsou jeho části propojeny
-- a vám pomůže promyslet architekturu \emph{před} psaním kódu.
Typické studentské projekty -- grafická hra, ovládání robota nebo zpracování dat
z fyzikálního pokusu -- mají vždy podobnou strukturu:
\textbf{vstupy}, \textbf{zpracování} a \textbf{výstupy}.
Právě tyto vztahy blokové schéma zachycuje.

\textbf{Jak sestavit blokové schéma:}
\begin{itemize}
	\item každý větší celek (modul, třída, soubor) = jeden blok
	\item šipky = tok dat nebo volání; na šipku napište \emph{co} teče nebo \emph{co} se volá
	\item datové soubory a herní stav odlište od modulů (jiný tvar nebo barva)
	\item schéma nemusí být úplné -- zachyťte hlavní vazby, ne každý řádek kódu
\end{itemize}

\textbf{Blokové schéma 1} -- procedurální program s datovými soubory (viz Příloha~B):

\begin{center}
\begin{tikzpicture}[
	node distance=1.4cm and 3cm,
	modul/.style={rectangle, draw, rounded corners, fill=blue!12,
	              minimum width=3.6cm, minimum height=1cm,
	              align=center, font=\small},
	data/.style={rectangle, draw, fill=orange!20,
	             minimum width=3.4cm, minimum height=0.8cm,
	             align=center, font=\small}
]
	\node[modul] (main)  {\texttt{main.c}\\hlavní program};
	\node[modul, right=of main] (lib) {\texttt{knihovna\_mereni.c}\\zpracování dat};
	\node[data,  below=of main] (out) {\texttt{mereni\_processed.csv}\\výsledky};
	\node[data,  below=of lib]  (in)  {\texttt{mereni.csv}\\surová data};

	\draw[->, bend left=15] (main) to
		node[above, font=\scriptsize] {volání funkcí} (lib);
	\draw[->, bend left=15] (lib)  to
		node[below, font=\scriptsize] {\texttt{Mereni[ ]}} (main);
	\draw[->] (main) -- node[left,  font=\scriptsize] {zapisuje} (out);
	\draw[->] (in)   -- node[right, font=\scriptsize] {čte}      (lib);
\end{tikzpicture}
\end{center}

\textbf{Blokové schéma 2} -- architektura jednoduché hry (stejný princip platí
pro robotické projekty: vstupy\,=\,senzory, zpracování\,=\,mikrokontrolér, výstupy\,=\,motory):

\begin{center}
\begin{tikzpicture}[
	node distance=1.4cm and 2.6cm,
	modul/.style={rectangle, draw, rounded corners, fill=blue!12,
	              minimum width=3.2cm, minimum height=1cm,
	              align=center, font=\small},
	data/.style={rectangle, draw, fill=orange!20,
	             minimum width=3.2cm, minimum height=0.8cm,
	             align=center, font=\small}
]
	\node[modul] (vstup)  {Vstup\\{\footnotesize klávesnice\,/\,myš}};
	\node[modul, right=of vstup]  (logika) {Herní logika\\{\footnotesize fyzika, kolize, AI}};
	\node[modul, right=of logika] (render) {Renderer\\{\footnotesize kreslení scény}};
	\node[data,  below=1.2cm of logika] (stav) {Herní stav\\{\footnotesize pozice, skóre, životy}};

	\draw[->]  (vstup)  -- node[above, font=\scriptsize] {události}     (logika);
	\draw[->]  (logika) -- node[above, font=\scriptsize] {data scény}   (render);
	\draw[<->] (logika) -- node[right, font=\scriptsize] {čte\,/\,zapisuje} (stav);
\end{tikzpicture}
\end{center}

Pro složitější schémata, která se obtížně kreslí v \LaTeX{}u, lze použít nástroj
\href{https://www.diagrams.net}{draw.io} a výsledek vložit jako obrázek.

\textbf{UML diagram tříd} se hodí pro objektově orientovaný kód.
Každá třída má tři části: název, atributy (proměnné) a metody (funkce).
Viditelnost se značí \texttt{+} (veřejné) nebo \texttt{-} (privátní).
Šipka s dutým trojúhelníkem vyjadřuje \emph{dědičnost} -- odvozená třída
přebírá atributy a metody rodičovské třídy a může je rozšířit nebo přepsat:

\begin{center}
\begin{tikzpicture}[font=\small]
	%% Rodičovská třída Entita
	\node[draw, thick, minimum width=4.4cm, align=center, font=\small\bfseries]
		(en) {Entita};
	\node[draw, thick, minimum width=4.4cm, below=-\pgflinewidth of en,
	      align=left, text width=4.0cm, inner sep=3pt]
		(ena) {
			\texttt{- x\,:\,int}\\
			\texttt{- y\,:\,int}
		};
	\node[draw, thick, minimum width=4.4cm, below=-\pgflinewidth of ena,
	      align=left, text width=4.0cm, inner sep=3pt]
		(enm) {
			\texttt{+ vykresli()}\\
			\texttt{+ pohyb()}
		};
	%% Odvozená třída Hrac
	\node[draw, thick, minimum width=4.4cm, align=center, font=\small\bfseries]
		at (-3.3, -5) (hn) {Hrac};
	\node[draw, thick, minimum width=4.4cm, below=-\pgflinewidth of hn,
	      align=left, text width=4.0cm, inner sep=3pt]
		(hna) {
			\texttt{- zivoty\,:\,int}\\
			\texttt{- skore\,:\,int}
		};
	\node[draw, thick, minimum width=4.4cm, below=-\pgflinewidth of hna,
	      align=left, text width=4.0cm, inner sep=3pt]
		(hnm) {
			\texttt{+ skoc()}\\
			\texttt{+ utoc()}
		};
	%% Odvozená třída Nepritel
	\node[draw, thick, minimum width=4.4cm, align=center, font=\small\bfseries]
		at (3.3, -5) (nn) {Nepritel};
	\node[draw, thick, minimum width=4.4cm, below=-\pgflinewidth of nn,
	      align=left, text width=4.0cm, inner sep=3pt]
		(nna) {
			\texttt{- rychlost\,:\,int}
		};
	\node[draw, thick, minimum width=4.4cm, below=-\pgflinewidth of nna,
	      align=left, text width=4.0cm, inner sep=3pt]
		(nnm) {
			\texttt{+ ai()}\\
			\texttt{+ utoc()}
		};
	%% Šipky dědičnosti (dutý trojúhelník = generalizace)
	\coordinate (spoj) at (0, -3.7);
	\draw (hn.north) -- ++(0, 0.5) -| (spoj);
	\draw (nn.north) -- ++(0, 0.5) -| (spoj);
	\draw[-{Triangle[open]}] (spoj) -- (enm.south);
\end{tikzpicture}
\end{center}

\textbf{Schéma relační databáze} má podobnou strukturu jako UML diagram tříd,
ale místo metod obsahuje datové typy sloupců.
Primární klíč (\textit{PK}) jednoznačně identifikuje každý záznam;
cizí klíč (\textit{FK}) odkazuje na primární klíč jiné tabulky.
Šipka s popiskem \textit{N\,:\,1} vyjadřuje, že jeden hráč může mít více rekordů:

\begin{center}
\begin{tikzpicture}[font=\small]
	%% Tabulka Hrac
	\node[draw, fill=green!15, minimum width=4.4cm, align=center,
	      font=\small\bfseries] (h0) {Hrac};
	\node[draw, minimum width=4.4cm, below=-\pgflinewidth of h0,
	      align=left, text width=4.0cm, inner sep=3pt]
		(h1) {\underline{\texttt{id}}\,:\,int\hfill\textit{PK}};
	\node[draw, minimum width=4.4cm, below=-\pgflinewidth of h1,
	      align=left, text width=4.0cm, inner sep=3pt]
		(h2) {\texttt{jmeno}\,:\,varchar};
	\node[draw, minimum width=4.4cm, below=-\pgflinewidth of h2,
	      align=left, text width=4.0cm, inner sep=3pt]
		(h3) {\texttt{heslo\_hash}\,:\,varchar};
	%% Tabulka Rekord
	\node[draw, fill=green!15, minimum width=4.4cm, align=center,
	      font=\small\bfseries] at (6.5, 0) (r0) {Rekord};
	\node[draw, minimum width=4.4cm, below=-\pgflinewidth of r0,
	      align=left, text width=4.0cm, inner sep=3pt]
		(r1) {\underline{\texttt{id}}\,:\,int\hfill\textit{PK}};
	\node[draw, minimum width=4.4cm, below=-\pgflinewidth of r1,
	      align=left, text width=4.0cm, inner sep=3pt]
		(r2) {\texttt{hrac\_id}\,:\,int\hfill\textit{FK}};
	\node[draw, minimum width=4.4cm, below=-\pgflinewidth of r2,
	      align=left, text width=4.0cm, inner sep=3pt]
		(r3) {\texttt{uroven}\,:\,int};
	\node[draw, minimum width=4.4cm, below=-\pgflinewidth of r3,
	      align=left, text width=4.0cm, inner sep=3pt]
		(r4) {\texttt{skore}\,:\,int};
	\node[draw, minimum width=4.4cm, below=-\pgflinewidth of r4,
	      align=left, text width=4.0cm, inner sep=3pt]
		(r5) {\texttt{datum}\,:\,date};
	%% Relace FK → PK  (trasa: nahoru pak doleva)
	\draw[->] (r2.west) |-
		node[pos=0.75, above, font=\scriptsize] {N\,:\,1} (h1.east);
\end{tikzpicture}
\end{center}
